% 18 November 2026. Prepared for Edward Baker with OpenAI Codex assistance.
% Actual model and effort unavailable; helper configuration inherited.
% E46 r1 integration of assessed Notes42,45,52; no new numerical execution.
\section{Correlated errors in complete-source discriminant enclosures}
\label{sec:correlated-enclosures}

The contour identities provide exact observables of one complete source.
Their useful accuracy depends on retaining that common source when errors
are propagated. Three distinct improvements are available: apply the actual
interpolation operator to nodal errors, combine the two moment readouts
before taking their rectangular support, and retain a common analytic
constraint on their annular aliases. We develop these analytical steps
from \cite{BakerCM42,BakerCM45,BakerCM52}; their internal assessments are
\cite{BakerCMReview37,BakerCMReview39,BakerCMReview43}. Numerical source
budgets and the inherited canonical example are supporting material in
Appendix~\ref{app:compact-example}.

These estimates refine a direct contour enclosure. They do not create an
additional observation or shrink the exact common-source feasible set in
Proposition~\ref{ts:sound}. A direct method given the same correlations
receives the same bounds. Throughout this section $m$ denotes an
interpolation degree and $n_\Gamma$ a contour node count; neither is a
finite-divisor cutoff.

\subsection{The exact interpolation operator and its error ledger}

Let $m\ge1$, let $T_k$ be the Chebyshev polynomial of the first kind, and
use the ascending Lobatto nodes and endpoint weights
\begin{equation}
 x_j=\cos\frac{(m-j)\pi}{m},\qquad
 d_0=d_m=\frac12,\quad d_j=1\ (0<j<m).
 \label{ce:nodes}
\end{equation}
The endpoint-weighted cosine transform defines the exact interpolant
$I_m y(x)=\sum_{k=0}^m a_kT_k(x)=\sum_{j=0}^m\ell_j(x)y_j$, where
\begin{align}
 a_k&=\frac{2d_k}{m}\sum_{j=0}^m d_jy_jT_k(x_j),\notag\\
 \ell_j(x)&=\frac{2d_j}{m}\sum_{k=0}^m d_kT_k(x_j)T_k(x).
 \label{ce:cardinals}
\end{align}
The weights of both the endpoint nodes and endpoint modes are present.

\begin{theorem}[A cardinal norm with an exact exterior value]
\label{ce:cardinal-norm}
The functions in \eqref{ce:cardinals} are the Lagrange cardinal
polynomials for \eqref{ce:nodes}. Their real operator norm satisfies
\begin{equation}
 \sum_{j=0}^m|\ell_j(x)|\le\sqrt{2m}\quad(-1\le x\le1),
 \qquad
 \sum_{j=0}^m|\ell_j(x)|=T_m(|x|)\quad(|x|\ge1).
 \label{ce:cardinal-bound}
\end{equation}
Consequently, for $X\ge1$ and complex errors $|e_j|\le E_{\rm node}$,
\begin{equation}
 \sup_{-X\le x\le1}|I_m e(x)|
 \le\Lambda_m(X)E_{\rm node},\qquad
 \Lambda_m(X):=\max\{\sqrt{2m},T_m(X)\}.
 \label{ce:nodal-propagation}
\end{equation}
The interior estimate need not be sharp.
\end{theorem}

\begin{proof}
The endpoint-half-weighted cosine sum is half the real sum over the
$2m$th roots of unity. The product-to-sum formula therefore gives
\begin{equation}
 \sum_{j=0}^m d_jT_k(x_j)T_l(x_j)
 =\delta_{kl}\frac{m}{2d_k},\qquad 0\le k,l\le m.
 \label{ce:orthogonality}
\end{equation}
Reversing the usual cosine grid multiplies its $k$th vector by
$(-1)^k$ and leaves this orthogonality unchanged. The square transform
is invertible, so \eqref{ce:cardinals} has
$\ell_j(x_i)=\delta_{ij}$. Applying \eqref{ce:orthogonality} once more
gives
\[
 \sum_{j=0}^m\frac{\ell_j(x)^2}{d_j}
 =\frac2m\sum_{k=0}^m d_kT_k(x)^2.
\]
Since $\sum d_j=m$ and $|T_k(x)|\le1$ on $[-1,1]$, weighted
Cauchy--Schwarz proves the first bound in \eqref{ce:cardinal-bound}.

For $x>1$, the Lagrange product has sign
$\operatorname{sgn}\ell_j(x)=(-1)^{m-j}=T_m(x_j)$.
For $x<-1$ its sign is $(-1)^j=(-1)^mT_m(x_j)$.
Exact interpolation of $T_m$ now gives respectively $T_m(x)$ and
$(-1)^mT_m(x)=T_m(|x|)$ for the sum of absolute values. Continuity
includes the endpoints. Finally $T_m(u)=\cosh(m\operatorname{arcosh}u)$
is increasing for $u\ge1$, and triangle inequality applies to arbitrary
complex $e_j$, proving \eqref{ce:nodal-propagation}.
\end{proof}

\begin{corollary}[The degree-$32$ continuation factor]
\label{ce:degree32}
On $[-503/500,1]$, the degree-$32$ operator in
\eqref{ce:cardinals} has norm strictly below $17$.
\end{corollary}
\begin{proof}
The interior bound is $8$. The recurrence
\[
 A_0=1,\quad A_1=503,\quad
 A_{k+1}=1006A_k-250000A_{k-1}
\]
gives $T_k(503/500)=A_k/500^k$. Its exact integer comparison
$A_{32}<17\cdot500^{32}$, recorded with numerator and denominator in
\cite{BakerCM45,BakerCMReview39}, proves the other part of the maximum
in \eqref{ce:nodal-propagation}. This is a finite rational inequality,
not an estimate inferred from a source sample.
\end{proof}

To apply this operator, fix a nonzero constant real multiple $G$ of $H$,
a real spatial reference $x_*$, and a real time map $t=t_c+b x$, $b>0$.
Let $P_j(w)$ be one exact spatial polynomial at each time node $t_j$.
Suppose the complete nodal error on a spatial disk is paid as
\begin{equation}
 |G(t_j,x_*+w)-P_j(w)|\le E_{\rm node},\qquad
 E_{\rm node}=E_{\rm sample}+E_{\rm alias}+E_{\rm spatial}.
 \label{ce:source-ledger}
\end{equation}
Here the first term includes source quadrature, its complete integration
tails and the input sample radii; the other two pay spatial aliasing and
spatial truncation. These meanings refer to the specified construction of
$P_j$, not to an arbitrary list of errors with the same names.
For $P(t,w)=I_m[P_j(w)](x)$, the exact decomposition
\[
 G(t,x_*+w)-P(t,w)
 =G(t,x_*+w)-I_m[G(t_j,x_*+w)](x)
   +I_m[G(t_j,x_*+w)-P_j(w)](x)
\]
and Theorem~\ref{ce:cardinal-norm} yield
\begin{equation}
 |G-P|\le E_{\rm time}+\Lambda_m(X)E_{\rm node}.
 \label{ce:complete-field-error}
\end{equation}

The true-field time remainder is a separate term. If
$|G(t,x_*+w)|\le M_R$ on $|w|\le R$ and on the convex hull of all time
nodes and evaluation times, then, for $|w|\le r<R$,
\begin{equation}
 E_{\rm time}\le
 \frac{b^{m+1}(2m+2)!M_R}{(m+1)!(R-r)^{2m+2}}
 \sup_{x\in[-X,1]}|\omega_m(x)|,
 \qquad
 \omega_m(x)=2^{1-m}(x^2-1)U_{m-1}(x).
 \label{ce:time-remainder}
\end{equation}
Indeed, $\omega_m=\prod_{j=0}^m(x-x_j)$, the heat equation gives
$\partial_t^{m+1}G=(-1)^{m+1}\partial_z^{2m+2}G$, and spatial
Cauchy estimates pay the displayed derivative. For complex $G$, project
the interpolation error at the chosen point onto its own complex
direction and apply the real scalar interpolation remainder; its
derivative is bounded by the same complex modulus. The time interval in
this argument includes any extrapolated real times. A positive-time
majorant alone does not pay a negative-time extension.

The exact polynomial in \eqref{ce:complete-field-error} must also be
distinguished from its arithmetic enclosure. Conversion of input numbers,
transform coefficients and polynomial evaluation must enclose that same
$P$. If coefficient balls and outward evaluation already contain those
errors, their radii are not added a second time to $E_{\rm node}$.
A midpoint-only implementation instead owes an explicit additional
arithmetic allowance. Neither the derivative weights of the polynomial
nor $E_{\rm time}$ are reduced merely by improving $\Lambda_m(X)$.
Differentiating a scalar error estimate requires an analytic remainder
bound on a larger spatial domain, as in the preceding approximation
sections. The quotient and branch-matched logarithm estimates in
Section~\ref{sec:compact} then convert the complete field ledger into
contour data. Every floor and count must hold on whole time-contour
cells before a downstream sign test.

\subsection{One logarithm and the sharp rectangular comparison}

Fix a time in an independently counted two-zero tube, put $w=z-x_*$,
and suppose the annulus $r_-\le|w|\le r_+$ contains no zeros, with
exactly the target pair inside its inner circle. Assume the same function
$\sigma G(t,x_*+w)/w^2$, for a fixed $\sigma\in\{1,-1\}$, has
strictly positive real part on both boundary circles. The harmonic
minimum principle extends this property to the annulus. Thus
\begin{equation}
 \mathcal L(w)=\Log\frac{\sigma G(t,x_*+w)}{w^2}-\kappa,
 \qquad \kappa\in\R,
 \qquad |\mathcal L|\le M_\pm\quad(|w|=r_\pm)
 \label{ce:annular-log}
\end{equation}
is one analytic centered logarithm, continuous on the boundaries. The
boundary constants must include the full field and arithmetic ledger;
they do not describe independently chosen analytic functions.

Let $r_-<\rho<r_+$ and define exact and discrete angular means by
\begin{equation}
 \mathsf T_0 f=\frac1{2\pi}\int_0^{2\pi}f(\rho\e^{\ii\theta})\dd\theta,
 \qquad
 \mathsf T_{n_\Gamma}f=
 \frac1{n_\Gamma}\sum_{j=0}^{n_\Gamma-1}f(w_j),
 \quad w_j=\rho\e^{2\pi\ii j/n_\Gamma}.
 \label{ce:means}
\end{equation}
Take the integer $n_\Gamma>2$. The target power sums relative to $x_*$
and their symmetric coordinates obey
\begin{equation}
 p_k=-k\,\mathsf T_0(w^k\mathcal L)\quad(k=1,2),
 \qquad \gamma=c-x_*=\frac{p_1}{2},\qquad
 D=\frac{p_2}{2}-\gamma^2.
 \label{ce:power-sums}
\end{equation}
To see the normalization, use $G_z/G=\mathcal L_z+2/w$ in the
argument-principle formula for $p_k$ and integrate the first term by
parts. The $2/w$ term integrates to zero for $k\ge1$.
Conjugation invariance makes these power sums real; no sign of $D$ is
assumed.

On a set $\mathcal J$ of conjugation representatives, write the enclosed
node values as
\begin{equation}
 \mathcal L(w_j)=\lambda_j+e_j,\qquad
 |\Re e_j|\le\alpha_j,\quad |\Im e_j|\le\beta_j.
 \label{ce:node-boxes}
\end{equation}
The symmetry weight $\nu_j$ is two for a nonreal conjugate pair and one
for a real node, so $\sum_{j\in\mathcal J}\nu_j=n_\Gamma$.
For a conjugation-compatible kernel $K$, define
\begin{equation}
 \mathcal S(K)=\frac1{n_\Gamma}\sum_{j\in\mathcal J}\nu_j
 \bigl(|\Re K(w_j)|\alpha_j+|\Im K(w_j)|\beta_j\bigr).
 \label{ce:box-support}
\end{equation}
This is the exact support of the absolute real kernel sum over the
rectangular relaxation \eqref{ce:node-boxes}: choosing each component
endpoint according to the sign of its real linear coefficient attains
the positive extreme, and reversing every choice attains its negative.
The true source errors need not fill this relaxation. Exactness of its
support is not physical independence of source perturbations.

Let $\epsilon_1,\epsilon_2$ be chosen positive annular quadrature
allowances for $p_1,p_2$. With the symmetry-weighted discrete means,
\begin{align}
 p_1&=-\Re\mathsf T_{n_\Gamma}(w\lambda)
      -\Re\mathsf T_{n_\Gamma}(we)+\eta_1,
 &|\eta_1|&\le\epsilon_1,\notag\\
 \frac{p_2}{2}&=-\Re\mathsf T_{n_\Gamma}(w^2\lambda)
      -\Re\mathsf T_{n_\Gamma}(w^2e)+\frac{\eta_2}{2},
 &|\eta_2|&\le\epsilon_2.
 \label{ce:moment-errors}
\end{align}
For the initial comparison the two $\eta$ intervals are independent
outer allowances. Put
\begin{align}
 \gamma_0&=-\tfrac12\Re\mathsf T_{n_\Gamma}(w\lambda),
 &s_{20}&=-\Re\mathsf T_{n_\Gamma}(w^2\lambda),\notag\\
 r_\gamma&=\tfrac12\{\mathcal S(w)+\epsilon_1\},
 &r_2&=\mathcal S(w^2)+\epsilon_2/2.
 \label{ce:marginals}
\end{align}
These give the sharp individual projections
$\Gamma=[\gamma_0-r_\gamma,\gamma_0+r_\gamma]$ and
$p_2/2\in[s_{20}-r_2,s_{20}+r_2]$ of that error relaxation.

\begin{proposition}[Sharp separate-moment baseline]
\label{ce:separate}
On the Cartesian product of those two marginal intervals, the exact
range of $D$ is obtained by subtracting the full real square range
$\Gamma^2$ from $[s_{20}-r_2,s_{20}+r_2]$.
If $\gamma_0>r_\gamma\ge0$, its endpoints and width are
\begin{align}
 L_{\rm sep}&=s_{20}-r_2-(\gamma_0+r_\gamma)^2,\notag\\
 U_{\rm sep}&=s_{20}+r_2-(\gamma_0-r_\gamma)^2,
 &W_{\rm sep}&=2r_2+4\gamma_0r_\gamma.
 \label{ce:separate-range}
\end{align}
\end{proposition}
\begin{proof}
The Cartesian product makes the second moment and $\gamma$ independently
selectable. The extrema therefore select the two endpoints of the
second-moment interval and the opposite extrema of $\gamma^2$.
On a positive $\Gamma$, squaring is increasing, giving
\eqref{ce:separate-range}. If $\Gamma$ crosses zero, its square has
minimum zero; the general square-range statement still applies.
\end{proof}

The separate projections can be individually sharp while their extreme
choices are mutually incompatible as images of the same node errors.
It is this particular loss of information, rather than gratuitous
rounding of a comparator, that the following centering addresses.

\subsection{A centered discriminant with its complete square}

Choose any fixed exact real $h$ and set
\begin{equation}
 J_h=\frac{p_2}{2}-hp_1+h^2=D+(\gamma-h)^2,
 \qquad
 J_0=s_{20}-2h\gamma_0+h^2.
 \label{ce:centered-identity}
\end{equation}
The shared node-error kernel is $w(w-h)$, not two independent
moment kernels. Define
\begin{equation}
 S=\mathcal S\bigl(w(w-h)\bigr),\qquad
 A=\epsilon_2/2+|h|\epsilon_1,
 \qquad [q_{\rm lo},q_{\rm hi}]=\operatorname{range}(\Gamma-h)^2.
 \label{ce:centered-budget}
\end{equation}

\begin{theorem}[Complete centered enclosure]
\label{ce:centered}
Under the preceding source, count and error hypotheses,
\begin{equation}
 D\in[J_0-S-A-q_{\rm hi},\ J_0+S+A-q_{\rm lo}].
 \label{ce:centered-range}
\end{equation}
If $h=\gamma_0$ and $\Gamma$ is exactly the symmetric interval in
\eqref{ce:marginals}, its width is $2(S+A)+r_\gamma^2$.
Under the same choice $h=\gamma_0$ and symmetric $\Gamma$, if
$\gamma_0>r_\gamma$, a sufficient strict improvement over
\eqref{ce:separate-range} is
\begin{equation}
 2\{r_2+2\gamma_0r_\gamma-S-A\}>r_\gamma^2.
 \label{ce:centered-gain}
\end{equation}
\end{theorem}
\begin{proof}
Substitute \eqref{ce:moment-errors} into
\eqref{ce:centered-identity}. The common node error is
$-\Re\mathsf T_{n_\Gamma}(w(w-h)e)$, bounded by $S$; the remaining
tail is $\eta_2/2-h\eta_1$, bounded by $A$. Subtract the entire real
range of $(\gamma-h)^2$ to obtain \eqref{ce:centered-range}.
For $h=\gamma_0$, that range is $[0,r_\gamma^2]$.
Subtracting the resulting width from \eqref{ce:separate-range} proves
\eqref{ce:centered-gain}.
\end{proof}

For an implemented reference choose an exact representable $h$, for
example the midpoint of an enclosing nominal-center interval. Its
displacement from $\gamma_0$ then remains in $\Gamma-h$; it is not
declared zero. Kernel, nominal and square arithmetic must all use that
same $h$ and be outward enclosed. If $h$ is instead supplied as an
uncertain interval, every occurrence must be bounded over its full
range. The square minimum is zero only when $0\in\Gamma-h$;
otherwise both endpoint magnitudes enter the minimum as well as the
maximum. A small square payment may not be discarded.

The theorem need not improve every input box. Moreover it bounds $J_h$
and its square remainder separately, so it does not claim the exact
range on the full joint source set. Applying the same centering twice
does not create further information. Direct extraction with the same
kernel is precisely this estimator.

\subsection{A joint annular alias estimate}

The preceding allowance $A$ still lets the two moment tails vary
independently. Analyticity of the common logarithm imposes a constraint
that survives without new source observations. Write its Laurent series
as $\mathcal L(w)=\sum_{k\in\mathbb Z}l_kw^k$ and put
\begin{equation}
 q_-=(r_-/\rho)^{n_\Gamma},\qquad
 q_+=(\rho/r_+)^{n_\Gamma},\qquad 0<q_\pm<1.
 \label{ce:alias-ratios}
\end{equation}

\begin{theorem}[Common-log Parseval bound]
\label{ce:joint-alias}
For $n_\Gamma>2$ and real $h$, the combined alias error satisfies
\begin{align}
 \left|(\mathsf T_{n_\Gamma}-\mathsf T_0)
                  \bigl((w^2-hw)\mathcal L\bigr)\right|
 &\le A_{\rm joint},\notag\\
 A_{\rm joint}
 &:=\sum_{\sigma\in\{-,+\}}M_\sigma r_\sigma
        \sqrt{r_\sigma^2+h^2}\,
        \frac{q_\sigma}{\sqrt{1-q_\sigma^2}}.
 \label{ce:joint-alias-bound}
\end{align}
In particular this bounds $|\eta_2/2-h\eta_1|$ in
\eqref{ce:moment-errors}, with its existing moment normalizations.
\end{theorem}
\begin{proof}
The discrete angular mean selects exactly the Laurent indices divisible
by $n_\Gamma$. Absolute convergence on the interior circle gives
\begin{equation}
 (\mathsf T_{n_\Gamma}-\mathsf T_0)
          \bigl((w^2-hw)\mathcal L\bigr)
 =\sum_{j\ne0}(l_{jn_\Gamma-2}-h l_{jn_\Gamma-1})
                    \rho^{jn_\Gamma}.
 \label{ce:alias-identity}
\end{equation}
The constant centering in \eqref{ce:annular-log} contributes no alias
because $n_\Gamma>2$. There is no extra factor two: the relevant
moments are $p_2/2$ and $hp_1$ in \eqref{ce:power-sums}.

On either boundary Parseval and \eqref{ce:annular-log} give
\begin{equation}
 \sum_{k\in\mathbb Z}|l_k|^2r^{2k}
 =\frac1{2\pi}\int_0^{2\pi}|\mathcal L(r\e^{\ii\theta})|^2\dd\theta
 \le M_r^2.
 \label{ce:parseval}
\end{equation}
For $j\ge1$ use $r=r_+$ in \eqref{ce:alias-identity}. The two
coefficient sequences in the normalized vector $(l_kr_+^k)_k$ have
weights $r_+^2q_+^j$ and $-hr_+q_+^j$. Their indices are disjoint:
an overlap would require a nonzero multiple of $n_\Gamma$ to equal
one. Their squared weight norm is therefore
\[
 (r_+^4+h^2r_+^2)\sum_{j\ge1}q_+^{2j}
 =(r_+^4+h^2r_+^2)\frac{q_+^2}{1-q_+^2}.
\]
For $j\le-1$, use the inner boundary and the analogous weights
$r_-^2q_-^{|j|}$ and $-hr_-q_-^{|j|}$. Cauchy--Schwarz on each
boundary and triangle inequality between the two contributions prove
\eqref{ce:joint-alias-bound}. Taking a real part or changing the common
sign only reduces the modulus.
\end{proof}

This proof retains a common boundary $L^2$ constraint across the aliased
coefficients. It does not make them independently realizable. The inner
and outer contributions are still bounded separately, so the result is
not asserted to be sharp for a single analytic function on the annulus.

For comparison, separate Cauchy coefficient bounds and geometric sums
give the unrounded allowance
\begin{equation}
 A_{\rm raw}=\sum_{\sigma\in\{-,+\}}M_\sigma
 r_\sigma(r_\sigma+|h|)\frac{q_\sigma}{1-q_\sigma}.
 \label{ce:raw-tail}
\end{equation}
When $\epsilon_1,\epsilon_2$ are upward enclosures of the corresponding
individual moment allowances, $A_{\rm raw}\le A$ in
\eqref{ce:centered-budget}. Stored allowances must not be identified
with exact physical errors or silently rounded downward.

\begin{corollary}[A uniform fractional tail reduction]
\label{ce:fractional-tail}
If $3/50\le |h|/r_\sigma\le1$ on both boundaries, then
\begin{equation}
 A_{\rm joint}<\frac{19}{20}A_{\rm raw}\le\frac{19}{20}A,
 \label{ce:tail-reduction}
\end{equation}
provided $A_{\rm raw}>0$.
\end{corollary}
\begin{proof}
For $x\in[0,1]$, the ratio $\sqrt{1+x^2}/(1+x)$ is decreasing.
For $x\ge3/50$ it is at most $\sqrt{2509}/53<19/20$, where the
strict inequality is the integer comparison $400\cdot2509<361\cdot2809$.
The ratio of the geometric factors is
$(1-q)/\sqrt{1-q^2}<1$ for $0<q<1$. Compare each nonzero positive
summand in \eqref{ce:joint-alias-bound} and \eqref{ce:raw-tail}.
\end{proof}

The ratio condition is part of this corollary, not a universal property
of a centered discriminant. Theorem~\ref{ce:joint-alias} remains valid
when it fails, including at $h=0$.

\subsection{Complete width and conservative endpoint translation}

\begin{theorem}[A retained-budget width improvement]
\label{ce:complete-width}
Keep the same $J_0,S,\Gamma,h$ and full square range in
\eqref{ce:centered-range}. Suppose the common annular information proves
$|\eta_2/2-h\eta_1|\le\vartheta A$ for a fixed $0\le\vartheta<1$.
Then replacing only $A$ by $\vartheta A$ gives a sound discriminant
enclosure whose exact displayed width is smaller by
\begin{equation}
 2(1-\vartheta)A.
 \label{ce:width-saving}
\end{equation}
Suppose $L^-,U^+$ are outward enclosures of the old displayed endpoints,
including all previous nominal, kernel and endpoint arithmetic, and
$0\le A_{\rm lo}\le A$ is a proved lower bound for the chosen old
allowance. Then the conservative translated interval
\begin{equation}
 \bigl[L^-+(1-\vartheta)A_{\rm lo},\quad
       U^+-(1-\vartheta)A_{\rm lo}\bigr]
 \label{ce:translated-endpoints}
\end{equation}
is also sound. Its width saving relative to $[L^-,U^+]$ is exactly
$2(1-\vartheta)A_{\rm lo}$ before any new display rounding.
\end{theorem}
\begin{proof}
The proof of Theorem~\ref{ce:centered} uses only a valid bound for the
combined tail; substitute $\vartheta A$. Each ideal endpoint moves
inward by $(1-\vartheta)A$, while the entire square range and node
support stay unchanged. This proves \eqref{ce:width-saving}.
The old lower endpoint enclosure is no larger than the ideal old lower
endpoint. Adding at most $(1-\vartheta)A$ therefore keeps it below the
ideal improved lower endpoint. The upper endpoint follows with the
opposite inequality. This proves \eqref{ce:translated-endpoints} without
discarding or reconstructing the old rounding slack.
\end{proof}

For serialized interval data, parse the old lower endpoint downward and
the old upper endpoint upward. Parse the positive allowance components
downward only to form $A_{\rm lo}$; these lower bounds are used for the
\emph{amount of justified inward movement}, not as error upper bounds.
Exact rational translation then preserves all inherited payments.
If a relative saving is claimed, compare a proved saving lower bound
with an upper bound for the old parsed width; a fresh rounded interval
requires its own outward rounding and comparison in the appropriate
directions. The condition $A_{\rm lo}>0$ gives strict inward movement.

The first-moment allowance inside $\Gamma$ is unchanged. In particular,
an improvement of the linear combination does not justify reducing
either moment marginal or its square range. Source integration tails,
spatial and time approximation, coefficient arithmetic, branch clearance,
uncertain coordinates and all whole-time estimates also remain paid.
An endpoint enclosure bounds a whole-path minimum only when an
independent whole-time premise, such as $D_t>0$ on the complete interval,
identifies that minimum with the endpoint. No derivative-width comparison
can replace this premise.

These theorems are generic analytic and error-set statements. A shifted
quadratic with a double target has $p_1=2b$, $p_2=2b^2$ and
$J_h=(b-h)^2$ at its event; retaining the full square returns $D=0$.
It can satisfy the ratio condition of Corollary~\ref{ce:fractional-tail}
and still collide. The canonical applications therefore retain the
separate source, count, branch and time hypotheses summarized in
Section~\ref{sec:discussion}. The estimates themselves do not distinguish
theta from every colliding control. Their established benefit is a
smaller specified enclosure with matched information, not a new
collision barrier, added height coverage, recursive contraction or a
cost-saving theorem. The inherited analytical width consequence and the
still-unexecuted rational endpoint transformation are distinguished in
the supporting appendix.
